% !TEX program = xelatex
\documentclass[10pt,a4paper]{article}

\usepackage[margin=0.65in]{geometry}
\usepackage{fontspec}
\usepackage{polyglossia}
\usepackage{enumitem}
\usepackage[hidelinks]{hyperref}
\usepackage{titlesec}

\setmainlanguage[numerals=maghrib]{arabic}
\setotherlanguage{english}

\newfontfamily\arabicfont[Script=Arabic,Scale=1.05]{Amiri}
\newfontfamily\arabicfontsf[Script=Arabic,Scale=1.05]{Amiri}
\newfontfamily\arabicfonttt[Script=Arabic,Scale=1.05]{Amiri}
\newfontfamily\englishfont{Latin Modern Roman}

\setlength{\parindent}{0pt}
\setlist[itemize]{leftmargin=1.2em, rightmargin=0pt, noitemsep, topsep=2pt}

\titleformat{\section}
  {\large\bfseries}
  {}
  {0em}
  {}
  [\titlerule]

\hypersetup{
    colorlinks=false,
    pdfborder={0 0 0}
}

\urlstyle{same}

\newcommand{\en}[1]{\textenglish{#1}}

\begin{document}

{\LARGE \textbf{عبدالرحمن السعيد الصاوي}}\\
مهندس روبوتات ومحاكاة مركبات\\
\en{ROS/ROS2} | \en{BeamNG.tech} | تكامل المحاكاة | جسور \en{API} | ضمان جودة آلي

البريد الإلكتروني: \href{mailto:abdelrahman.elsaid.elsawy@gmail.com}{\en{abdelrahman.elsaid.elsawy@gmail.com}}\\
الهاتف/\en{WhatsApp}: \href{https://wa.me/905365231899}{\en{+90 536 523 1899}}\\
الموقع: \href{https://abdelrahmanelsaidelsawy.github.io/index.html}{\en{abdelrahmanelsaidelsawy.github.io}}\\
\en{LinkedIn}: \href{https://www.linkedin.com/in/abdelrahman-elsaid-elsawy}{\en{linkedin.com/in/abdelrahman-elsaid-elsawy}}\\
\en{GitHub}: \href{https://github.com/AbdelrahmanElsaidElsawy}{\en{github.com/AbdelrahmanElsaidElsawy}}

\section*{الملخص}

مهندس روبوتات وبرمجيات بخبرة تزيد عن 6 سنوات في \en{ROS/ROS2} وتكامل أنظمة الروبوتات والأنظمة الذاتية ومحاكاة المركبات عالية الدقة. خبرة في تكامل \en{BeamNG.tech} وسير عمل \en{MATLAB/Simulink} وتطوير جسور \en{API} والمحاكاة اللحظية واكتساب البيانات واختبار ضمان الجودة الآلي وتصحيح الأخطاء والتوثيق التقني والدعم الفني للعملاء. مهارات في \en{C/C++} و\en{Python} و\en{Lua} و\en{Bash} و\en{Git} و\en{Docker} و\en{Ubuntu} و\en{Gazebo} و\en{RViz} و\en{MATLAB} و\en{Simulink}.

\section*{المهارات الأساسية}

\textbf{الروبوتات:} \en{ROS}، \en{ROS2}، \en{SLAM}، \en{MoveIt}، \en{Gazebo}، \en{RViz}، الملاحة الذاتية، دمج الحساسات\\
\textbf{المحاكاة:} \en{BeamNG.tech}، محاكاة المركبات، محاكاة عالية الدقة، سير عمل المحاكاة\\
\textbf{البرمجة:} \en{C}، \en{C++}، \en{Python}، \en{Lua}، \en{SQL}، \en{Bash}\\
\textbf{التكامل:} جسور \en{API}، تكامل الأنظمة، محاكاة لحظية، اكتساب بيانات، \en{MATLAB}، \en{Simulink}\\
\textbf{الاختبار وضمان الجودة:} اختبار ضمان جودة آلي، اختبارات انحدار، إعادة إنتاج الأخطاء، تصحيح الأخطاء، تحليل الأسباب الجذرية\\
\textbf{الأدوات:} \en{Git}، \en{Docker}، \en{Ubuntu}، \en{Zendesk}، التوثيق التقني

\section*{الخبرة المهنية}

\textbf{\en{Research Software Engineer}} \hfill نوفمبر 2021 -- حتى الآن\\
\textbf{\en{BeamNG GmbH}} \hfill عن بُعد

\begin{itemize}
    \item دمج \en{BeamNG.tech} مع \en{ROS/ROS2} و\en{MATLAB} و\en{Simulink} لمحاكاة مركبات عالية الدقة وسير عمل روبوتات خارجية.
    \item تطوير وصيانة جسور \en{API} للتحكم اللحظي بالمحاكاة واكتساب البيانات والتفاعل مع المركبات وتكامل الأدوات الخارجية.
    \item تنفيذ اختبار ضمان جودة آلي وتحليل انحدار وتحقيق في المشكلات باستخدام \en{C++} و\en{Python} و\en{Lua}.
    \item التحقيق في المشكلات المبلَّغ عنها من العملاء والمكتشفة داخلياً عبر إعادة إنتاج الأخطاء وتحليل الأسباب الجذرية وتقديم إصلاحات أو حلول مؤقتة أو إرشاد تقني.
    \item دعم مستخدمي \en{B2B} العالميين عبر \en{Zendesk} والمنتديات التقنية حول المحاكاة وواجهات البرمجة والتكاملات وسير العمل.
    \item تحسين توثيق \en{API} والدروس والأمثلة الإرشادية وإرشادات \en{MATLAB/Simulink} والمواد التقنية للمطورين.
    \item تحويل مشكلات ضمان الجودة والدعم المتكررة إلى حالات اختبار قابلة لإعادة الاستخدام وتحسينات توثيق وملاحظات منتج قابلة للتنفيذ.
    \item المساهمة في مستودعات هندسة \en{BeamNG.tech} بما في ذلك \en{BeamNGpy} وتوثيق \en{API} وسير عمل المحاكاة وإصلاحات ضمان الجودة والدروس وأمثلة التكامل.
    \item تنفيذ والتحقق من سير عمل تحميل السيناريوهات والمهام، بما في ذلك اكتشاف السيناريو وتحميله واسترجاع المركبات وربط الحساسات وربط المركبة والتحكم النصي واستطلاع الحساسات.
\end{itemize}

\textbf{\en{Mechatronics Systems Engineer}} \hfill فبراير 2020 -- يوليو 2021\\
\textbf{\en{Ideasparq Robotics Sdn Bhd}} \hfill في الموقع

\begin{itemize}
    \item تطوير أنظمة تحكم قائمة على \en{ROS} وحلول ملاحة ذاتية لروبوتات الخدمة.
    \item تنفيذ خوارزميات دمج الحساسات لتطبيقات الروبوتات.
    \item اختبار منصات روبوتية باستخدام \en{Gazebo} و\en{RViz} و\en{Docker} و\en{Ubuntu}.
    \item تكامل الأنظمة عبر البرمجيات والحساسات ومنطق التحكم والمحاكاة وبيئات النشر.
\end{itemize}

\textbf{\en{Autonomous Robotics Engineer}} \hfill يونيو 2019 -- يناير 2020\\
\textbf{\en{ENDSTRUCT}} \hfill في الموقع

\begin{itemize}
    \item تطوير أنظمة ملاحة ذاتية لروبوتات متنقلة ضمن مشروع \en{GOVTECH}.
    \item تطوير أنظمة مراقبة الحالة الصحية لمنصات روبوتية متنقلة.
\end{itemize}

\textbf{مساعد باحث} \hfill سبتمبر 2019 -- ديسمبر 2019\\
\textbf{\en{Autonomous Vehicle Laboratory}} \hfill \en{UMPSA Universiti Malaysia Pahang Al-Sultan Abdullah}

\begin{itemize}
    \item معالجة بيانات \en{LiDAR} ثلاثية الأبعاد للكشف اللحظي عن الأجسام.
    \item تطوير وسيط برمجي قائم على \en{ROS} لتجريد الحساسات.
\end{itemize}

\section*{مساهمات تقنية مختارة}

\begin{itemize}
    \item تكامل \en{BeamNG.tech} مع \en{ROS/ROS2} و\en{MATLAB} و\en{Simulink}.
    \item تطوير \en{BeamNGpy} وجسور \en{API} للتحكم اللحظي بالمحاكاة واكتساب البيانات.
    \item التحقق من سير عمل تحميل السيناريوهات والمهام للمحاكاة الآلية.
    \item ضمان جودة آلي واختبارات انحدار وإعادة إنتاج الأخطاء وتصحيحها وتحليل الأسباب الجذرية.
    \item دعم تقني لمستخدمي \en{B2B} العالميين عبر \en{Zendesk} والمنتديات التقنية.
    \item توثيق وأمثلة تكامل لواجهات البرمجة وسير عمل المحاكاة و\en{MATLAB/Simulink} ومواد \en{HIL/DIL}.
    \item تكامل أنظمة روبوتات باستخدام \en{ROS} و\en{Gazebo} و\en{RViz} و\en{Docker} و\en{Ubuntu}.
\end{itemize}

\section*{التعليم}

\textbf{بكالوريوس هندسة الميكاترونكس والروبوتات والأتمتة}
\hfill 2016 -- 2020\\
\en{UMPSA Universiti Malaysia Pahang Al-Sultan Abdullah}

\section*{منشور علمي}

\textbf{\en{The Flexural Strength Prediction of Porous Cu-Sn-Ti Composites via Artificial Neural Networks}}\\
\en{The 6th International Conference on Robotics Intelligence and Applications, 2018}\\
\en{DOI:} \href{https://doi.org/10.1007/978-981-13-8323-6_34}{\en{10.1007/978-981-13-8323-6\_34}}

\section*{دورات وورش عمل}

\begin{itemize}
    \item \en{Learn ROS2 as a ROS1 Developer and Migrate Your ROS Projects -- Udemy, Nov 2022}
    \item \en{Level 2: Simulation and Hardware Implementation for SLAM -- Udemy, Nov 2022}
    \item \en{Getting Started with Azure DevOps Boards -- Coursera, Oct 2022}
    \item \en{Command Line in Linux -- Coursera, Sep 2022}
    \item \en{Arduino Workshop -- Core Electronics, Sep 2022}
    \item \en{Open-Source Software Development Methods -- The Linux Foundation, Oct 2021}
    \item \en{Python Scripting for DevOps Specialization -- LearnQuest, Sep 2021}
    \item \en{Python Scripting: Files, Inheritance, and Databases -- LearnQuest, Sep 2021}
    \item \en{Web Application Technologies and Django -- University of Michigan, Sep 2021}
    \item \en{AWS Cloud Technical Essentials -- AWS, Aug 2021}
    \item \en{DevOps and Build Automation with Python -- LearnQuest, Aug 2021}
    \item \en{Python Scripting: Dates, Classes, and Collections -- LearnQuest, Aug 2021}
    \item \en{Configuration Management and the Cloud -- Google, Aug 2021}
    \item \en{Introduction to Git and GitHub -- Google, Jul 2021}
    \item \en{Technical Support Fundamentals -- Google, Jul 2021}
    \item \en{Troubleshooting and Debugging Techniques -- Google, Jul 2021}
    \item \en{Using Python to Interact with the Operating System -- Google, Jun 2021}
    \item \en{Crash Course on Python -- Google, May 2021}
    \item \en{Malaysia Autonomous Vehicle Challenge (MAVeC'19) -- MARii, Nov 2018}
    \item \en{The International Conference on Robotics Intelligence and Applications (RiTA 2018) -- KIAST SG, 2018}
    \item \en{Boundless Innovation Week -- HPMM, Aug 2018}
    \item \en{CAM Workshop -- The Institution of Engineers, Malaysia, Nov 2017}
    \item \en{IEM MATLAB Workshop -- The Institution of Engineers, Malaysia, Oct 2017}
    \item \en{Arduino for Beginners Workshop -- SPARKTECH Team at IIUM, Nov 2016}
    \item \en{Online Database Class: SCOPUS -- University Malaysia Pahang, Nov 2016}
\end{itemize}

\section*{اللغات}

العربية: لغة أم أو إتقان ثنائي اللغة\\
الإنجليزية: إتقان مهني للعمل\\
التركية: مستوى أساسي\\
الألمانية: مستوى أساسي\\
الماليزية: إتقان محدود للعمل
\end{document}
